%! Author = sriadityavedantam
%! Date = 3/30/26

\section{Subgroups}

\begin{definition}[Subgroup]
    Let $G$ be a group and let $\varnothing\neq H\subseteq G$.
    We say $H$ is a subgroup of $G$ if $H$ is a group under the operation of $G$.
    That is, for all $a,b\in H$, we have $ab\in H$, $1\in H$, and if $a\in H$, then $a^{-1}\in H$.
    Associativity is inherited from $G$.
    This is denoted by
    \[
        H\leq G.
    \]
\end{definition}

Something to note that you may realize through a lot of practice problems is that
\[
    H\leq G
\]
if and only if for all $a,b\in H$,
\[
    ab^{-1}\in H.
\]
We could also have a subgroup $H\leq G$ with $\ord(H)<\infty$.
Then $H\leq G$ if and only if $H\neq\varnothing$ and $H$ is closed under the operation of the group.

\begin{example}
    Let $H$ be a finite nonempty subset of a group $G$ that is closed under the group operation.
    Then $H$ is a subgroup of $G$.
    Indeed, let $a\in H$.
    Since $H$ is finite, the sequence
    \[
        a,a^2,a^3,\dots
    \]
    must repeat, so $a^m=a^n$ for some $m>n$.
    Then
    \[
        a^{m-n}=1.
    \]
    Hence
    \[
        a^{-1}=a^{m-n-1}\in H.
    \]
    So every element of $H$ has an inverse in $H$, and therefore $H\leq G$.
\end{example}

$G=GL(n,\mathbb{F})$ is the group of nonsingular $n\times n$ matrices with entries in $\mathbb{F}$.
This is called the General Linear Group.
One of our exercises is that
\[
    SL(n,\mathbb{F})
\]
is the set of matrices in $GL(n,\mathbb{F})$ with determinant $1$.
This is called the Special Linear Group.
Since
\[
    \det(AB)=\det(A)\det(B),
\]
if $\det(A)=\det(B)=1$, then
\[
    \det(AB)=1.
\]
Also, if $\det(A)=1$, then
\[
    \det(A^{-1})=\frac{1}{\det(A)}=1.
\]
So $SL(n,\mathbb{F})\leq GL(n,\mathbb{F})$.
Let $\mathbb{C}^*$ be the set of nonzero complex numbers.
Let $n\in\mathbb{Z}^{>0}$.
Then
\[
    \left\{
        \exp\left(\frac{2k\pi i}{n}\right):k=0,\dots,n-1
    \right\}
\]
is the set of complex $n$th roots of unity.
Also,
\[
    \exp\left(\frac{2k\pi i}{n}\right)
    =
    \cos\left(\frac{2k\pi}{n}\right)
    +
    i\sin\left(\frac{2k\pi}{n}\right).
\]
We can think of multiplication here as cyclical.
Given any group, there are important examples you can get from any group.

\begin{definition}[Cyclic Subgroup Generated by $a$]
    Let $a\in G$.
    The cyclic subgroup generated by $a$, denoted by $\gen{a}$, is
    \[
        \gen{a}\coloneqq \{a^n:n\in\mathbb{Z}\}.
    \]
\end{definition}

We need $n\in\mathbb{Z}$ since we want all powers, positive and negative, because the order can be finite or infinite.
Since it is closed under multiplication and inverses, it is also a subgroup.
We consider this as the smallest subgroup of $G$ containing $a$.
We call this cyclic because when the order is finite, the powers repeat in a cycle.
For an infinite subgroup, look at $(\mathbb{Z},+)$.
The subgroup generated by $2$ is
\[
    \dots,-4,-2,0,2,4,\dots.
\]
For example, if $H\leq G$ and $a\in H$, then
\[
    \gen{a}\leq H.
\]
\begin{definition}[Subgroup Generated by a Subset]
    If $S\subseteq G$, the subgroup generated by $S$, denoted by $\gen{S}$, is the smallest subgroup of $G$ containing $S$.
    This is the same as
    \[
        \bigcap\limits_{H\leq G \text{ and } S\subseteq H} H,
    \]
    which is the intersection of all subgroups of $G$ containing $S$.
\end{definition}

For example, if we take two appropriate elements of a symmetric group, they may generate the entire group.
Similar to ideals, we also have trivial subgroups.
For example, ponder looking at the subgroup generated by the identity element.

\begin{definition}
    If $G$ is a group and
    \[
        G=\gen{a}
    \]
    for some $a\in G$, then $G$ is cyclic.
\end{definition}

$(\mathbb{Z}_n,+)$ is cyclic, since it is generated by $1$.
It is also generated by any $m\in\mathbb{Z}_n$ such that
\[
    \gcd(m,n)=1.
\]
Also,
\[
    \mathbb{Z}_5^*=\{1,2,3,4\}
\]
is cyclic, generated by $2$ or $3$.
For
\[
    \mathbb{Z}_7^*=\{1,2,3,4,5,6\},
\]
we note that
\[
    2^3=8\equiv 1\mod 7,
\]
so $2$ does not generate the whole group.
However,
\[
    3^2=9\equiv 2\mod 7,
\]
and since $2^3\equiv 1\mod 7$, we get
\[
    3^6\equiv 1\mod 7.
\]
In fact, the powers of $3$ run through all nonzero classes mod $7$, so $\ord(3)=6$.
Since
\[
    \ord(\mathbb{Z}_7^*)=6,
\]
the group is cyclic.

\begin{theorem}{
    If $p$ is prime, then $\mathbb{Z}_p^*$ is cyclic.
}
    This is a standard theorem in group theory and number theory.
    The proof is more extensive than what we need right now, but the result is important to know.
\end{theorem}